\subsection{Kolej 17.~listopadu}

Celý komplex budovy koleje a menzy 17. listopadu je tvořen dvěma bloky koleje (budovy A a B) a budovou dnes již
bývalé menzy (budova C, nyní patří FHS a probíhá její přestavba). Pod všemi
budovami se nachází dvě patra rozsáhlého suterénu, kde je umístěna menza,
technické zázemí koleje, sklady prádla i jiného, tělocvična, hudební zkušebna,
kotelna, dílny a spousta dalších různě využitých či zcela nevyužitých prostor

Budovy A a B jsou v suterénu spojeny (aktuálně ne vždy)průchozí cestou, zpočátku hůře
nalezitelnou. 

\subsubsection{Budova A}
Budova A je větší z obou budov (20 pater + přízemí), blíže k cestě. Doprava je
zajišťována čtyřmi výtahy (většinou jezdí jenom tři, v době největší potřeby i
méně) a dvěma požárními schodišti potaženými linem. Na koleji existuje nepsané pravidlo (psané doporučení)
nejezdit výtahem z přízemí do 3. (a nižšího) patra, nepoužívat výtah na přejezd
do sousedních pater a analogicky nejezdit z 5. a nižšího patra do přízemí, také
jezdit do -1 mimo doby oběda není slušné. Kromě toho, že tyto cesty jsou po
schodech opravdu rychlejší, jejich dodržováním si ušetříte vražedné pohledy od
zbytku výtahu, který obvykle jede do dvouciferných pater. Samozřejmě se
připouští výjimky, ať už ze zdravotních důvodů, fyzického vyčerpání či podobně
opodstatněného důvodu.

Okna pokojů jsou orientována buď na východ, s malebným výhledem na magistrálu,
nebo na západ s výhledem na Céčko a tunel Blanka (resp. na kolonu před
portálem). Východní strana je pro velkou část obyvatel populárnější, jelikož v
letních měsících pálí do oken pouze ranní slunce.

Na dvacátém patře v budově A se tradičně, ač většinou bez
povolení, pořádají tak dvakrát do roka chodbovice – párty na chodbě. Aby byli
všichni ušetřeni starostí s oprávněnými stížnostmi, raději se tam vůbec
nestěhujte, pokud víte, že by vám to vadilo.

V přízemí je lab.
Více informací o labu najdete na stránce o labech. Za výtahy se nachází
kanceláře.

Jedna smutnější informace zde musí také zaznít, ač se tónem nehodí do veselého
zbytku kuchařky. V roce 2008 ukončil svůj život student skokem z 19. patra
budovy A, obdobně skončil v roce 2016 i život studenta ze 16. patra. Nelze k
tomu dodat nic jiného než – nedělejte to.


\subsubsection{Budova B}
Budova B je ta menší z budov (16 pater + přízemí), kdysi bývala synonymem vyšší
životní úrovně, ale to už nějakou dobu neplatí.

V podzemí je místnost Kolejní rady, klavír, posilovna a cestou k Áčku hudební 
zkušebna. V přízemí pak sídlí kolejní obchůdek, studovna, ekumenická místnost
 a jeden ze sedmi vedoucích správy KaM (podorgán ředitele KaM).

Doprava je zajištěna dvěma malými výtahy a výtahem evakuačním. Relativně často
jeden z nich nejezdí do všech pater (hlavně 15. patro) anebo nejezdí vůbec.
 Je tu také požární schodiště, pro velký úspěch opět potaženo linem. Oproti
výtahům na budově A, kde se můžete setkat s cedulkou
„MIMO PROVOZ“, se zde můžete setkat s cedulkou „VÝTAH V PROVOZU“. Skrytý význam
této cedulky ponecháme pozornému čtenáři za domácí cvičení.

Pokoje mají výhled buď na východ, tedy na Miladu (rozpadlá budova, která už dávno nemá stát), nebo na západ do stráně.
Oproti západní straně budovy A jsou nižší patra západní strany kryta strání před
ostrým sluncem.

\subsubsection{Bývalá menza (budova C)}
Nyní je to staveniště. Aktuálně studentům neslouží, ale měla by se sem (po dokončení
stavby) přesunout Fakulta humanitních studií (taktéž FHS neboli Fakulta
hledající smysl), která dosud vlastní budovu nemá. Vzhledem k tomu, že jsou
všechny budovy propojené, je hluk ze staveniště intenzivní i v pokojích, které
výhled na staveniště nemají.


\subsubsection{Magistrála}
Magistrála vede z Jižního Města na Prosek. Je opravdovým požehnáním naší koleje. Ničí naše
uši a otravuje naše plíce. Je použitelná na mnoho způsobů. V zimních měsících
nám poskytuje rozptýlení, když sledujeme boj silničářů o udržení sjízdnosti této
důležité komunikace. Ve zkouškovém období a při výuce cizích jazyků se bavíme
počítáním projíždějících aut. Navíc poskytuje blahodárný nízkofrekvenční zvuk
přispívající ke klidnému usínání. Její osvětlení v noci připomíná řecké písmeno
epsilon.

Skoro každé ráno na magistrále vzniká fronta spěchajících řidičů, pohybující se
rychlostí pomalého chodce. Na psychiku některých pak působí příznivě, mohou-li
zdeptané řidiče pěšky předcházet.

K magistrále se po nekonečných letech stavby a průtahů připojil i tunelový
komplex Blanka. V odpoledních hodinách před jeho vjezdem a za jeho výjezdem
obvykle vzniká kolona.


\subsubsection{Technické speciality}

Zlí jazykové tvrdí, že otevře-li se daný nadkritický počet oken nad sebou, kolej
spadne. Zatím to však nebylo experimentálně ověřeno. Sama budova má podobné
vlastnosti jako komín. Zatímco v nejnižších patrech bývala zima, v osmnáctém
patře nebylo dokonce ani v neizolované budově za největších mrazů vůbec potřeba
otevírat topení (problémy byly pochopitelně opačné, topení obvykle nešlo
zavřít). 

Meteorologové vám jistě rádi vysvětlí, proč celá kolej funguje jako větrná
hůrka. Zatímco v celé Praze fouká jen mírný větřík, kolem Áčka stojí matfyzáci,
kteří se snaží v silných poryvech větru alespoň udržet se na nohou. Ti
šťastnější občas udělají i malý krok dopředu. Celá cesta kolem (kratší!) stěny
Áčka trvá až několik minut, vydrží jen ti nejodolnější. Po několika větrných
dnech i nejtvrdohlavější matfyzák pochopí, proč není vždy dobré otevírat okna na
ventilaci bez ověření, že nebude třeba zapotřebí celé buňky k jeho zavření.

Stavební práce na koleji (či v její blízkosti) všeobecně začínají se zkouškovým
obdobím; těžko říci zda náhodou či úmyslem. Ať už to bylo budování
protipovodňových opatření v zimě 2009, zmíněná přestavba vestibulů v létě 2009,
pokládání základů pro budovu pod kolejí v létě 2014, stavba tunelu Blanka anebo
přestavba budovy C.


\subsubsection{Internet}

V současnosti je ubytovaným k dispozici symetrické připojení 100 (Mb) na 100
(Mb) za 100 (Kč) s veřejnou IP adresou. Celá kolej pak má desetigigabitové
připojení do PASNETu (agregace je tedy cca 1:10). Tomu celému se většinou říká
prostě \textit{KolejNet}. 

Další info ohledně připojení ke kabelovému unternetu najdete na:\\
\url{www.matfyzak.cz/kucharka?title=Kolej_17._listopadu#Internet}

Nově je ale možnost na všechno ohledně kabelového internetu zapomenout a připojit se na koleji na wifi - eduroam, 
což je wifi jednotná pro všechny studenty UK. Pozor, pokud přestanete být studentem, 
(i např. než se později zapíšete na magisterské studium) tak vám dočasně eduroam seberou. 


\subsubsection{Eduroam}
Připojení na eduroam je podstatně snažší. Na Albeři byste měli dostat heslo k Casu,
kde si mužete poté nastavit heslo na eduroam. Dále již vše funguje, jako když se přihlašujete na jakoukoli wifi.
POZOR!! přihlašovací jméno je ve tvaru: jméno@cuni.cz nebo čísloisicu@cuni.cz


\subsubsection{Opravy závad}
Opravy na naší koleji, kupodivu, probíhají docela rychle. Nejdůležitější je
zapsat poruchu na vrátnici na lísteček. Opravná akce probíhá asi takto: ráno
zapíšete do sešitu tekoucí sifon u umyvadla. Kolem poledne se na pokoji objeví
zámečník a ujistí vás, že na tohle je potřeba instalatér. Areál koleje je velký,
takže než je instalatér nalezen, má po směně. Druhý den ráno se dostaví
instalatér. Zkonstatuje, že záchod je v pořádku a zanechá vám o tom vzkaz.
Můžete tedy znovu připomenout na vrátnici, že stále protéká sifon. Takto lze
závadu obecně vyřešit, protože instalatérovi po nějaké době dojdou výmluvy a
sifon opraví. Obecně platí, že máte-li jakoukoli závadu (včetně maličkostí typu
ucpaný odpad, nefunkční klíč od skříně či zavzdušněné topení), napište to do
závad. Až na závady vyžadující výměnu – může se stát, že díly nejsou na skladě,
a tedy si budete muset počkat, dokud se neobjednají – je problém obvykle vyřešen
do druhého dne.

\textbf{Hlavně nic neopravujte sami!}


\subsubsection{Kolejní rada}
Kolejní radu, svůj samosprávný orgán, který zastupuje zájmy ubytovaných proti
zájmům vedení koleje, si volí na koleji ubytovaní studenti. Na jejích webových
stránkách najdete spoustu užitečných věcí – telefonní seznam, formuláře,
provozní řády, výtahovou etiketu, otevírací doby, vyhledávání ubytovaných a
simulátor rozestavování nábytku. Adresa: \url{listopad.koleje.cuni.cz}.


\subsubsection{Zamlouvání pokojů}
Celá akce probíhá přes Internet někdy v srpnu. Nejdříve dostanou šanci ti,
kteří jsou spokojeni se svým příbytkem a nechtějí se v příštím roce stěhovat jinam.
Poté si ti, kteří se stěhovat chtějí, vyberou z toho, co zbylo, systémem, kdo
dřív přijde, ten si lépe vybere. Obzvláště pokoje na východní straně mizí s

překvapivou rychlostí a vězte, že v 5. minutě je již téměř každá buňka na této
straně obsazena alespoň jedním člověkem. 

Takto to fungovalo do roku 2018, kdy bylo nutno přistoupit k obnově systému. Od té doby to nefunguje a zámluvy tedy probíhají přes email kolem poloviny prázdnin s ubytovací kanceláří.

Celé se to netýká studentů, kteří na koleji bydlí i přes prázdniny, neboť
ti si
mohou pokoj, na kterém právě bydlí, zamluvit již v červnu při podpisu smlouvy na
celý další rok.

Z důvodu příchodu GDPR (a lenosti KaM něco programovat) jsou všichni studenti
ubytováni do stejných pokojů jako v předchozím roce. Noví obyvatelé jsou do
pokojů přiděleni náhodně.


\subsubsection{Co si lze půjčit}
U vrátného si můžete zapůjčit klíče od prádelny či sušárny (praček je řádově
méně než sušáren, navíc pračky se postupem času rozbíjejí), stejně jako vysavač,
žehličku a pumpičku na kolo. Na prádelny a sušárny bývá fronta, občas je lepší
předem na vrátnici zavolat.

Pro uschování jízdních kol slouží tzv. kolárny, které spravuje kolejní rada.
Jsou skoro na každém druhém patře, ale některé pořád plné. Kontakty na správce
najdete na stránkách kolejní rady.

Také máme na koleji knihovnu sci-fi a fantasy literatury. Nachází se ve dvacátém
patře Áčka a je otevřená každé pondělí a středu od 19:00 do 20:00. Více
informací najdete na jejích stránkách knihovna.krakonos.org, nebo na nástěnce
kolejní rady u požárního výtahu na Béčku.


\subsubsection{Svoboda na naší koleji}
Je značná, téměř už taková, jakou bychom chtěli. Smí se (přesněji, je
tolerováno) libovolně si přestavět pokoj a vyzdobit si ho dle libosti. Na
základě vyhlášky magistrátu o čistotě hlavního města Prahy je však zakázána
výzdoba oken. Rovněž na chodby a do výtahů je zakázáno cokoliv vylepovat. Taktéž
je zakázáno lepit cokoliv, nebo, nedej bože, zatloukat hřebíčky do stěn či
dveří; ale koho by ty holé zdi bavily...

Dřívější měsíční poplatek za elektrospotřebiče byl s velkou slávou zahrnut do
ceny kolejného, přirozeně s patřičnou přirážkou. S běžnými spotřebiči tedy
nebývá problém, jen na lednici potřebujete potvrzení lékaře.

Chov zvířat je střídavě zakazován a povolován hygieniky. Většinou je třeba
souhlas alespoň spolubydlícího, občas celé buňky. A pak samozřejmě vedení
koleje. Informujte se v ubytovací kanceláři.


\subsubsection{Sportování}
Z důvodu přerodu ve staveniště se v areálu koleje již nenachází hřiště na
volejbal nebo nohejbal, dva basketbalové koše mezi Béčkem a Céčkem ani jakési
pseudohřiště na fotbal. Zkusit ale můžete tenisové kurty u řeky, kde ale platíte
stejně jako ostatní (s kolejí mají společnou jenom polohu). Letité pokusy o
změnu situace příliš ovoce nepřinesly, což je způsobeno zejména tím, že v okolí
koleje není zcela jasné, co komu patří. Pokud pak hledáte nářadí, např. na
volejbal, pak na vrátnici budovy B najdete věci k zapůjčení.

Pravidelně nefungující a nespolehlivé výtahy inspirovaly studenty k uspořádání
několika forem běhu do schodů. Sportovněji zaměřený (a originálně pojmenovaný)
je „Běh do schodů“, při kterém se běží od menzy do 20. patra na čas (poslední
rekord z podzemí do 20. patra je minuta a 36 vteřin). Druhou takovouto akcí je
„KolejCup“, pravidelně pořádaný během zkouškového, který je vědomostní obdobou,
kdy jsou po koleji rozmístěny otázky a na základě (ne)znalosti se mezi nimi týmy
přesouvají. Obsahem otázek je především učivo z prvního ročníku napříč všemi
obory. 

\subsubsection{Kuchařka}
Kolej 17. listopadu má svoji vlastní podrobnou kuchařku v režii Kolejní rady.
Najdete ji na adrese \url{listopad.koleje.cuni.cz/kucharka}.